\scanpage{38}\mspage{25}

\S\,2. \uline{Description de} $C_I\subset\mathrm{Fl}(M(I^{\unsure{o|\circ}}))$, \uline{pour} $I$ \uline{ensemble ordonné}.

Soit $M$ une catégorie, stable par (petites) $\unsure{\underrightarrow{\lim}|\underline{\lim}}$, $C\subset\mathrm{Fl}\,M$ \inserted{(``cofibrations'').} On suppose

\[
\left\{\begin{aligned}
 &\text{a) } C\text{ contient les iso et stable par \unsure{cochgmt|cochangement} de base}\\
 &\text{b) } C\text{ stable par composition transfinie.}
\end{aligned}\right. \tag*{(2.1.)}
\]

Soit $I$ un \inserted{(petit)} ensemble ordonné, et considérons

\[
 M(I^o)=\underline{\mathrm{Hom}}\,(I;\deleted{\unsure{C}}M),
\]

on va définir une partie

\[
 C_I\subset\mathrm{Fl}\,(M(I^o))
\]

ainsi, pour $u_\bullet:X_\bullet\longrightarrow Y$ flèche \unsure{de|dans} $M(I^o)$

\[
\begin{aligned}
 u_\bullet\in C_I\ &\Longleftrightarrow\ \forall\,i_0\in I\ \text{et pour tt ouvert}\\
 &J\subset I_{/i_0}=K,\ \text{la flèche}\\
 &u_{J,\deleted{\unsure{I_{/i_0}}}K}\ \text{est}\ \in C
\end{aligned}\tag*{(2.2.)}
\]

Voici la définition de $u_{J,K}$. Pour tt ouvert $J$ de $I$, on pose

\[
\left\{\begin{aligned}
 X_J&=\underset{i\in J}{\underrightarrow{\lim}}\,X_{\unsure{i|\bullet}},\quad
 Y_J=\underset{i\in J}{\underrightarrow{\lim}}\,Y_i,\\[-2pt]
 &\deleted{\text{et on considère la}}\\[-2pt]
 u_J&:X_J\longrightarrow \deleted{\unsure{X}}Y_J,\quad
 u_J=\underset{i\in J}{\underrightarrow{\lim}}\,u_i,
\end{aligned}\right.\tag*{(2.3.)}
\]

et on considère, pour $\deleted{\unsure{J}}K\subset\deleted{\unsure{K}}J$, le diagramme

% archive sheet number: 37

